%%
%% paper76_mandelbrot_mlc_de.tex
%% Autor: Michael Fuhrmann
%% Projekt: Specialist Mathematics Project
%% Thema: Die Mandelbrot-Menge — Lokale Zusammenhängigkeit und die MLC-Vermutung
%% Sprache: Deutsch
%% Erstellt: 2026-03-12
%% Letzte Änderung: 2026-03-12
%% Build: 164
%%
%% Beschreibung:
%%   Deutschsprachige Version von Paper 76. Behandelt die Mandelbrot-Menge M
%%   im Parameterraum der quadratischen Familie f_c(z) = z^2 + c, ihre globale
%%   Zusammenhängigkeit (Douady-Hubbard 1982), die MLC-Vermutung (M ist lokal
%%   zusammenhängend), Yoccoz' Satz für endlich-renormalisierbare Parameter,
%%   Thurston's topologisches Modell und offene Probleme.
%%

\documentclass[12pt,a4paper]{amsart}

%% Pakete für Mathematik, Layout, Encoding und Links
\usepackage{amsmath,amssymb,amsthm,mathtools,enumitem,booktabs,array,hyperref}
\usepackage[utf8]{inputenc}
\usepackage[T1]{fontenc}
\usepackage[ngerman]{babel}
\usepackage{geometry}
\geometry{margin=2.5cm}

%% Theorem-Umgebungen (englische Basis für Nummerierung)
\newtheorem{theorem}{Theorem}[section]
\newtheorem{lemma}[theorem]{Lemma}
\newtheorem{corollary}[theorem]{Korollar}
\newtheorem{proposition}[theorem]{Proposition}

%% Deutsche Theorem-Umgebungen
\newtheorem{satz}[theorem]{Satz}
\newtheorem{korollar}[theorem]{Korollar}
\newtheorem{vermutung}[theorem]{Vermutung}

\theoremstyle{definition}
\newtheorem{definition}[theorem]{Definition}
\newtheorem{definition_de}[theorem]{Definition}
\newtheorem{beispiel}[theorem]{Beispiel}

\theoremstyle{remark}
\newtheorem{bemerkung}[theorem]{Bemerkung}
\newtheorem{remark}[theorem]{Bemerkung}
\newtheorem{conjecture}[theorem]{Vermutung}

%% Mathematische Operatoren
\DeclareMathOperator{\diam}{diam}
\DeclareMathOperator{\dist}{dist}
\DeclareMathOperator{\Mod}{mod}

%% Hyperref-Konfiguration
\hypersetup{
    colorlinks=true,
    linkcolor=blue,
    citecolor=blue,
    urlcolor=blue,
    pdftitle={Die Mandelbrot-Menge: Lokale Zusammenhängigkeit und die MLC-Vermutung},
    pdfauthor={Michael Fuhrmann}
}

%% -------------------------------------------------------------------------
%% Abstract VOR \maketitle
%% -------------------------------------------------------------------------
\begin{abstract}
Wir geben eine umfassende Darstellung der Mandelbrot-Menge $\mathcal{M}$,
dem kanonischen Parameterraum der quadratischen Familie $f_c(z) = z^2 + c$.
Nach der Einführung gefüllter Julia-Mengen und der Charakterisierung von
$\mathcal{M}$ durch das Nichtentweichen des kritischen Orbits präsentieren
wir den Beweis von Douady und Hubbard (1982), dass $\mathcal{M}$ zusammenhängend
ist, mittels der Böttcher-Koordinate und der konformen Abbildung
$\Phi_{\mathcal{M}} \colon \widehat{\mathbb{C}} \setminus \mathcal{M}
\to \widehat{\mathbb{C}} \setminus \overline{\mathbb{D}}$.
Wir formulieren die \emph{MLC-Vermutung} — $\mathcal{M}$ ist lokal
zusammenhängend — und erläutern ihre weitreichenden Konsequenzen:
die Dichtheit hyperbolischer Abbildungen in der quadratischen Familie
und Thurston's topologische Charakterisierung von $\mathcal{M}$
durch Laminationen. Yoccoz' Satz (MLC für höchstens endlich-renormalisierbare
Parameter), die Fatou-Shishikura-Ungleichung und die Struktur der
hyperbolischen Komponenten werden ausführlich behandelt.
Wir schließen mit einer Übersicht offener Probleme.
\end{abstract}

\title{Die Mandelbrot-Menge: Lokale Zusammenhängigkeit\\
und die MLC-Vermutung}
\author{Michael Fuhrmann}
\address{Specialist Mathematics Project, Build 165}
\email{specialist-maths@project.local}
\date{2026-03-12}

\maketitle

%% -------------------------------------------------------------------------
\section{Einleitung}
%% -------------------------------------------------------------------------

Die \emph{Mandelbrot-Menge} $\mathcal{M}$, eingeführt von Benoit Mandelbrot
(1980)~\cite{Mandelbrot1980} und rigoros untersucht von Douady und Hubbard
ab 1982~\cite{DouadyHubbard1982}, ist das bekannteste Objekt der holomorphen
Dynamik. Sie parametrisiert die quadratische Familie
\[
    f_c(z) = z^2 + c, \quad c \in \mathbb{C},
\]
indem sie diejenigen Parameter markiert, für die das dynamische Verhalten
von $f_c$ „beschränkt" bleibt. Obwohl ihre Definition elementar ist,
kodiert $\mathcal{M}$ eine außerordentliche Fülle komplex-analytischer
Struktur. Ihre globale Zusammenhängigkeit wurde im ersten Paper von Douady
und Hubbard bewiesen, doch eine der grundlegendsten Fragen — ob $\mathcal{M}$
\emph{lokal} zusammenhängend ist — ist nach mehr als vier Jahrzehnten
offen geblieben.

Diese \emph{MLC-Vermutung} (``$\mathcal{M}$ is Locally Connected'') ist nicht
bloß eine topologische Kuriosität. Sie ist äquivalent zur \emph{Dichtheit der
Hyperbolizität} in der quadratischen Familie: der Aussage, dass jedes
$c \in \mathcal{M}$ durch hyperbolische Parameter approximiert werden kann.
Darüber hinaus würde MLC eine vollständige topologische Charakterisierung
von $\mathcal{M}$ via Thurston's Laminationstheorie liefern.

%% -------------------------------------------------------------------------
\section{Die Quadratische Familie und Julia-Mengen}
%% -------------------------------------------------------------------------

\subsection{Grundlegende Definitionen}

\begin{definition_de}[Gefüllte Julia-Menge]
Für $c \in \mathbb{C}$ ist die \emph{gefüllte Julia-Menge} von $f_c$:
\[
    K_c = \{ z \in \mathbb{C} \mid \sup_{n \geq 0} |f_c^n(z)| < \infty \},
\]
und die \emph{Julia-Menge} ist $J_c = \partial K_c$.
\end{definition_de}

Da $f_c$ ein Polynom vom Grad 2 ist, ist das Komplement
$\mathbb{C} \setminus K_c$ das \emph{Anziehungsbecken des Unendlichen}
$\mathcal{A}_\infty(c)$. Die Julia-Menge $J_c$ ist die Grenze zwischen
beschränkten und unbeschränkten Orbits.

\begin{satz}[Fatou-Julia-Dichotomie]
Für jedes $c \in \mathbb{C}$ gilt entweder:
\begin{enumerate}[label=\roman*)]
\item $K_c$ ist zusammenhängend und $J_c$ ist eine zusammenhängende,
      perfekte, nirgends dichte Teilmenge von $\mathbb{C}$, oder
\item $K_c$ ist eine Cantor-Menge (total unzusammenhängend).
\end{enumerate}
\end{satz}

\begin{proof}[Beweisskizze]
Die Böttcher-Koordinate liefert eine konforme Abbildung des Anziehungsbeckens
des Unendlichen auf $\mathbb{C} \setminus \overline{\mathbb{D}}$. Falls $0 \in K_c$
(der kritische Punkt entweicht nicht), ist $K_c$ voll und zusammenhängend.
Falls $0 \notin K_c$, zeigt die globale Böttcher-Abbildung, dass $K_c$ eine
Cantor-Menge ist (unter Verwendung von Sullivan's Satz über wandernde Gebiete).
\end{proof}

\subsection{Die Mandelbrot-Menge als Parameterraum}

\begin{definition_de}[Mandelbrot-Menge]
\[
    \mathcal{M} = \{ c \in \mathbb{C} \mid f_c^n(0) \not\to \infty \}
              = \{ c \in \mathbb{C} \mid K_c \text{ ist zusammenhängend} \}.
\]
\end{definition_de}

\begin{proposition}
$\mathcal{M}$ ist kompakt und $\mathcal{M} \subset \overline{D}(0, 2)$.
\end{proposition}

\begin{proof}
Für $|c| > 2$ gilt $|f_c(0)| = |c| > 2$, und eine Standardabschätzung
zeigt $|f_c^n(0)| \to \infty$. Die Kompaktheit folgt aus der Stetigkeit
von $c \mapsto f_c^n(0)$ und der Definition als Durchschnitt abgeschlossener
Mengen.
\end{proof}

%% -------------------------------------------------------------------------
\section{Die Böttcher-Koordinate und Green-Funktion}
%% -------------------------------------------------------------------------

\begin{satz}[Böttcher-Koordinate]
Für jedes $c \in \mathbb{C}$ existiert eine eindeutige konforme Abbildung
\[
    \phi_c \colon \mathcal{A}_\infty(c) \to \mathbb{C} \setminus \overline{\mathbb{D}}
\]
mit $\phi_c(f_c(z)) = \phi_c(z)^2$ und $\phi_c(z)/z \to 1$ für $z \to \infty$.
\end{satz}

Die \emph{Green-Funktion} von $K_c$ ist
\[
    G_c(z) = \log|\phi_c(z)| = \lim_{n\to\infty} \frac{1}{2^n} \log|f_c^n(z)|.
\]

\begin{definition_de}[Externe Strahlen]
Für $\theta \in \mathbb{R}/\mathbb{Z}$ ist der \emph{externe Strahl des Winkels
$\theta$} für $K_c$:
\[
    \mathcal{R}_c(\theta) = \phi_c^{-1}\bigl(\{ r e^{2\pi i \theta} \mid r > 1 \}\bigr).
\]
Ein externer Strahl \emph{landet} bei $\gamma_c(\theta) \in J_c$, falls
$\phi_c^{-1}(r e^{2\pi i \theta}) \to \gamma_c(\theta)$ für $r \to 1^+$.
\end{definition_de}

Die Dynamik von $f_c$ auf Strahlen ist besonders einfach: $f_c$ bildet
den Strahl des Winkels $\theta$ auf den Strahl des Winkels $2\theta \pmod 1$ ab.

%% -------------------------------------------------------------------------
\section{Zusammenhängigkeit der Mandelbrot-Menge}
%% -------------------------------------------------------------------------

\subsection{Die konforme Abbildung von Douady-Hubbard}

\begin{satz}[Douady-Hubbard 1982~\cite{DouadyHubbard1982}]
\label{satz:DH}
Es existiert eine Biholomorphie
\[
    \Phi_{\mathcal{M}} \colon \widehat{\mathbb{C}} \setminus \mathcal{M}
    \longrightarrow \widehat{\mathbb{C}} \setminus \overline{\mathbb{D}}
\]
mit $\Phi_{\mathcal{M}}(c)/c \to 1$ für $c \to \infty$, explizit gegeben
durch $\Phi_{\mathcal{M}}(c) = \phi_c(c)$.
\end{satz}

\begin{proof}[Beweisskizze]
Für $c \notin \mathcal{M}$ liegt der kritische Wert $c$ selbst im
Anziehungsbecken des Unendlichen von $f_c$, sodass $\phi_c(c)$ wohldefiniert
ist. Man verifiziert, dass $c \mapsto \phi_c(c)$ holomorph in $c$ auf
$\mathbb{C} \setminus \mathcal{M}$ ist (Impliziter Funktionssatz auf die
Böttcher-Funktionalgleichung angewandt), injektiv ist (ein sorgfältiges
Argument, das die Eindeutigkeit des kritischen Punkts $0$ benutzt), und
das korrekte asymptotische Verhalten hat. Die Erweiterung auf die Ränder
folgt aus dem Satz von Carathéodory.
\end{proof}

\begin{korollar}[Zusammenhängigkeit von $\mathcal{M}$]
Die Mandelbrot-Menge $\mathcal{M}$ ist zusammenhängend.
\end{korollar}

\begin{proof}
Die Existenz einer konformen Abbildung vom Komplement von $\mathcal{M}$
auf das Komplement von $\overline{\mathbb{D}}$ impliziert nach Riemann's
Theorie einfach zusammenhängender Gebiete, dass $\mathcal{M}$ zusammenhängend
und voll ist.
\end{proof}

\subsection{Externe Strahlen der Mandelbrot-Menge}

Die Abbildung $\Phi_{\mathcal{M}}$ definiert externe Strahlen für $\mathcal{M}$:
\[
    \mathcal{R}_{\mathcal{M}}(\theta) = \Phi_{\mathcal{M}}^{-1}\bigl(\{ r e^{2\pi i \theta} \mid r > 1 \}\bigr).
\]

\begin{satz}[Landung rationaler Strahlen]
Jeder externe Strahl $\mathcal{R}_{\mathcal{M}}(\theta)$ mit
$\theta \in \mathbb{Q}/\mathbb{Z}$ landet bei einem Misiurewicz-Punkt
oder an der Wurzel einer hyperbolischen Komponente.
\end{satz}

%% -------------------------------------------------------------------------
\section{Hyperbolische Komponenten und innere Struktur}
%% -------------------------------------------------------------------------

\begin{definition_de}[Hyperbolische Abbildung]
$f_c$ ist \emph{hyperbolisch}, falls jeder kritische Orbit zu einem
anziehenden Zyklus konvergiert. Die \emph{Hauptkardioid} $\mathcal{H}_0$
und der \emph{Periode-2-Knospe} $\mathcal{H}_2$ sind die größten
hyperbolischen Komponenten von $\mathcal{M}$.
\end{definition_de}

\begin{satz}[Parametrisierung hyperbolischer Komponenten]
Jede hyperbolische Komponente $\mathcal{H}$ von $\mathcal{M}$ ist einfach
zusammenhängend, und es existiert eine Biholomorphie
$\lambda_{\mathcal{H}} \colon \mathcal{H} \to \mathbb{D}$,
die \emph{Multiplikatorkarte}, die $c$ auf den Multiplikator $\lambda$
des anziehenden Zyklus von $f_c$ abbildet.
\end{satz}

\begin{proof}
Für die Hauptkardioid $\mathcal{H}_0$ hat $f_c$ einen anziehenden Fixpunkt
$z_c$ mit Multiplikator $\lambda = f_c'(z_c) = 2z_c$. Aus $z^2+c=z$ folgt
$z_c = (1 - \sqrt{1-4c})/2$, sodass die Multiplikatorkarte durch
$c = \lambda/2 - \lambda^2/4$ gegeben ist, eine Biholomorphie von
$\mathcal{H}_0$ nach $\mathbb{D}$. Der allgemeine Fall folgt durch dasselbe
Argument mit dem Impliziten Funktionssatz.
\end{proof}

%% -------------------------------------------------------------------------
\section{Die MLC-Vermutung}
%% -------------------------------------------------------------------------

\subsection{Formulierung und äquivalente Aussagen}

\begin{vermutung}[MLC — Mandelbrot Lokal Zusammenhängend]
\label{verm:MLC}
Die Mandelbrot-Menge $\mathcal{M}$ ist lokal zusammenhängend: Für jedes
$c_0 \in \mathcal{M}$ und jede Umgebung $U$ von $c_0$ existiert eine
zusammenhängende Umgebung $V \subseteq U$ von $c_0$ mit
$V \cap \mathcal{M}$ zusammenhängend.
\end{vermutung}

Die lokale Zusammenhängigkeit von $\mathcal{M}$ ist äquivalent zu:

\begin{satz}[MLC-Äquivalenzen]
Folgende Aussagen sind äquivalent:
\begin{enumerate}[label=\roman*)]
\item $\mathcal{M}$ ist lokal zusammenhängend (MLC).
\item Die konforme Abbildung $\Phi_{\mathcal{M}}^{-1}$ setzt sich stetig
      bis zum Rand fort: $\partial \mathbb{D} \to \partial \mathcal{M}$
      stetig und surjektiv.
\item Jeder externe Strahl von $\mathcal{M}$ landet, und die Landungsabbildung
      $\theta \mapsto \gamma_{\mathcal{M}}(\theta)$ ist stetig.
\item Das kombinatorische Modell $\mathcal{M}_{\mathrm{komb}}$
      (Thurston's abstraktes Mandelbrot-Modell) ist homöomorph zu $\mathcal{M}$.
\end{enumerate}
\end{satz}

\begin{proof}[Beweisskizze]
Die Äquivalenz von (i)--(iii) ist klassische Funktionentheorie:
für eine kompakte zusammenhängende Menge $K$ mit zusammenhängendem
Komplement ist lokale Zusammenhängigkeit äquivalent zur Carathéodory-
Erweiterungseigenschaft. Die Äquivalenz mit (iv) ist der Inhalt von
Thurston's Theorie topologischer Polynome~\cite{Thurston1985}.
\end{proof}

\subsection{Dichtheit der Hyperbolizität}

\begin{vermutung}[Dichtheit der Hyperbolizität]
\label{verm:Hyp}
Hyperbolische Parameter sind dicht in $\mathcal{M}$: Für jedes
$c_0 \in \mathcal{M}$ und jedes $\varepsilon > 0$ existiert ein
$c \in \mathcal{M}$ mit $|c - c_0| < \varepsilon$ und $f_c$ hyperbolisch.
\end{vermutung}

\begin{satz}[MLC impliziert Dichtheit der Hyperbolizität]
Vermutung~\ref{verm:MLC} impliziert Vermutung~\ref{verm:Hyp}.
\end{satz}

\begin{proof}[Beweisskizze]
Unter MLC ist der Eindruck jedes externen Strahls ein Punkt. Die Ränder
hyperbolischer Komponenten sind dicht in $\partial\mathcal{M}$ (Folge der
Parametrisierung und des Landungssatzes für rationale Strahlen). Da $\mathcal{M}$
unter MLC keine „wandernden" Fatou-Komponenten besitzt (Sullivan's Satz),
ist jede Nichtperiodizitätsmenge nirgends dicht, was die Dichtheit
hyperbolischer Parameter liefert.
\end{proof}

%% -------------------------------------------------------------------------
\section{Thurston's Topologisches Modell}
%% -------------------------------------------------------------------------

\subsection{Laminationen}

\begin{definition_de}[Quadratisch invariante Lamination]
Eine \emph{quadratisch invariante Lamination} (QIL) ist eine abgeschlossene
Kollektion $\mathcal{L}$ von Sehnen von $\overline{\mathbb{D}}$
(``Blätter'') mit:
\begin{enumerate}[label=\roman*)]
\item Blätter kreuzen sich nicht in $\mathbb{D}$.
\item $\mathcal{L}$ ist invariant unter $z \mapsto z^2$ auf $\partial \mathbb{D}$.
\item Das Vorwärtsbild jedes Blatts ist ein Blatt.
\end{enumerate}
Die \emph{Mandelbrot-Lamination} $\mathcal{L}_{\mathcal{M}}$ ist diejenige
QIL, die die Kombinatorik von $\mathcal{M}$ kodiert.
\end{definition_de}

\begin{satz}[Thurston~\cite{Thurston1985}]
Der Quotient $\partial \mathbb{D} / \sim_{\mathcal{L}_{\mathcal{M}}}$
ist ein topologisches Modell $\mathcal{M}_{\mathrm{komb}}$, das unter
MLC homöomorph zu $\mathcal{M}$ ist.
\end{satz}

%% -------------------------------------------------------------------------
\section{Yoccoz' Satz}
%% -------------------------------------------------------------------------

Das bedeutendste Teilergebnis in Richtung MLC stammt von
Jean-Christophe Yoccoz (1990):

\begin{satz}[Yoccoz~\cite{HubbardSchleicher1994}]
\label{satz:Yoccoz}
Sei $c_0 \in \mathcal{M}$ ein Parameter, der weder im Abschluss einer
hyperbolischen Komponente liegt noch unendlich renormalisierbar ist.
Dann ist $\mathcal{M}$ bei $c_0$ lokal zusammenhängend.
\end{satz}

\begin{proof}[Beweisidee]
Yoccoz führt das \emph{Yoccoz-Puzzle} ein: eine Folge verschachtelter
topologischer Scheiben $P^0 \supset P^1 \supset \cdots$ in der dynamischen
Ebene von $f_{c_0}$, entlang externen Strahlen zu Urbildern des
$\alpha$-Fixpunkts geschnitten. Die \emph{Yoccoz-Ungleichung}
\[
    \Mod(P^{n+1} \setminus \overline{P^{n+2}}) \geq
    \Mod(P^n \setminus \overline{P^{n+1}})
\]
zeigt, dass die Moduli nach unten beschränkt sind, was erzwingt, dass die
Puzzle-Stücke auf einen Punkt schrumpfen. Die entsprechenden Stücke im
\emph{Parameterraum-Puzzle} schrumpfen ebenfalls, was MLC bei $c_0$ beweist.
Die Endlichkeit der Renormalisierung ist genau die Hypothese, die ein
Degenerieren des Puzzles verhindert.
\end{proof}

\begin{korollar}
MLC gilt für alle Misiurewicz-Punkte und alle $c_0$ mit nicht-rekurrentem
kritischem Punkt ($0 \notin \omega(0)$ unter $f_{c_0}$).
\end{korollar}

%% -------------------------------------------------------------------------
\section{Renormalisierung und unendlich renormalisierbare Parameter}
%% -------------------------------------------------------------------------

\begin{definition_de}[Renormalisierbare Abbildung]
$f_c$ ist \emph{$n$-renormalisierbar}, falls eine topologische Scheibe
$U \ni 0$ existiert, sodass $f_c^n \colon U \to f_c^n(U)$ eine
polynomähnliche Abbildung vom Grad 2 mit zusammenhängender Julia-Menge ist.
$f_c$ ist \emph{unendlich renormalisierbar}, falls es für unendlich viele $n$
$n$-renormalisierbar ist.
\end{definition_de}

Yoccoz' Satz lässt den Fall unendlich renormalisierbarer Parameter offen.
Für diese ist kein allgemeiner Beweis von MLC bekannt. Teilergebnisse:

\begin{satz}[Lyubich 1997~\cite{Lyubich1997}]
Für Parameter von \emph{beschränktem kombinatorischem Typ} (d.h.\ die
Renormalisierungsperioden sind beschränkt) gilt MLC.
\end{satz}

\begin{satz}[Kahn-Lyubich 2009~\cite{KahnLyubich2009}]
Die Julia-Menge $J_c$ ist für jedes höchstens endlich renormalisierbare
$c \in \mathcal{M}$ lokal zusammenhängend.
\end{satz}

Der unendlich renormalisierbare Feigenbaum-Parameter
$c_{\mathrm{Feig}} \approx -1{,}4011552$ ist das primäre Hindernis:
MLC bei diesem Parameter würde aus Schranken an den konformen Moduli
des Feigenbaum-Fixpunkts folgen.

\begin{satz}[Dudko--Lyubich, 2023~\cite{DudkoLyubich2023}]
\label{satz:DudkoLyubich}
Die Mandelbrot-Menge $\mathcal{M}$ ist lokal zusammenhängend beim
Feigenbaum-Parameter $c = c_{\mathrm{Feig}}$.
\end{satz}

\begin{proof}[Beweisidee]
Dudko und Lyubich beweisen ausreichend starke Schranken an den konformen
Moduli des \emph{Feigenbaum-Fixpunkts} — des Fixpunkts des
Periodenverdopplungs-Renormalisierungsoperators.  Konkret zeigen sie,
dass das Hauptnest der Yoccoz-Puzzle bei $c_{\mathrm{Feig}}$ Moduli
hat, die gegen Unendlich wachsen; nach Lyubichs Kriterium von 1997
impliziert dies lokale Zusammenhängigkeit.  Entscheidend ist eine
rigorose \emph{computergestützte Schätzung} des Feigenbaum-Fixpunkts,
aufbauend auf Arbeiten von Lanford (1982) und der Sullivan--McMullen
Starrheitstheorie.
\end{proof}

\begin{bemerkung}
Satz~\ref{satz:DudkoLyubich} ist ein bedeutendes Teilergebnis: MLC ist
nun beim Feigenbaum-Punkt bewiesen.  Der allgemeine Fall für beliebige
unendlich renormalisierbare Parameter mit \emph{unbeschränkter Kombinatorik}
bleibt offen.  Der Preprint erschien im September 2023 auf arXiv
(arXiv:2309.02107).
\end{bemerkung}

%% -------------------------------------------------------------------------
\section{Die Fatou-Shishikura-Ungleichung}
%% -------------------------------------------------------------------------

\begin{satz}[Fatou-Shishikura-Ungleichung~\cite{Shishikura1987}]
Für jede rationale Abbildung $f$ vom Grad $d \geq 2$ gilt
\[
    \#\{\text{nicht-abstoßende Zyklen}\} \leq 2d - 2.
\]
\end{satz}

\begin{proof}[Beweis für Grad 2]
Für $f_c$ mit $d=2$ gilt $2d-2=2$ kritische Punkte (mit Vielfachheit).
Shishikura's Beweis verwendet quasikonforme Deformationstheorie: jeder
nicht-abstoßende Zyklus „absorbiert" mindestens einen kritischen Orbit
durch ein holomorphes Bewegungsargument, und verschiedene nicht-abstoßende
Zyklen absorbieren verschiedene kritische Orbits. Für $f_c$ ergibt dies
höchstens einen endlichen nicht-abstoßenden Zyklus.
\end{proof}

\begin{korollar}
Jedes $f_c$ hat höchstens einen anziehenden oder indifferenten Zyklus.
\end{korollar}

%% -------------------------------------------------------------------------
\section{Strukturresultate über $\partial\mathcal{M}$}
%% -------------------------------------------------------------------------

\subsection{Misiurewicz-Punkte}

\begin{definition_de}[Misiurewicz-Punkt]
$c_0 \in \partial \mathcal{M}$ ist ein \emph{Misiurewicz-Punkt}, falls
$0$ unter $f_{c_0}$ strikt präperiodisch ist: $f_{c_0}^{m+k}(0) = f_{c_0}^m(0)$
für gewisse $m \geq 1$, $k \geq 1$.
\end{definition_de}

\begin{satz}
Misiurewicz-Punkte sind dicht in $\partial \mathcal{M}$.
\end{satz}

\subsection{Hausdorff-Dimension von $\partial\mathcal{M}$}

\begin{satz}[Shishikura 1998~\cite{Shishikura1998}]
Die Hausdorff-Dimension des Randes der Mandelbrot-Menge ist 2:
\[
    \dim_H(\partial \mathcal{M}) = 2.
\]
\end{satz}

\subsection{Parabolische Implosion}

An parabolischen Parametern (wo $f_c$ einen parabolischen Zyklus mit
Multiplikator $e^{2\pi i p/q}$ besitzt) erfährt die Julia-Menge $J_c$
eine dramatische Veränderung: unendlich viele winzige Kopien von $J_c$
implodieren in das Fatou-Gebiet, ein Phänomen bekannt als
\emph{parabolische Implosion}~\cite{LavaursRousseau1993}. Dies erklärt
die charakteristische „Antennen"-Struktur in der Nähe parabolischer
Punkte in $\partial\mathcal{M}$.

%% -------------------------------------------------------------------------
\section{Offene Probleme}
%% -------------------------------------------------------------------------

Wir sammeln die wichtigsten offenen Probleme in diesem Themengebiet.

\begin{vermutung}[MLC, Douady~\cite{DouadyHubbard1982}]
\label{verm:MLC2}
Die Mandelbrot-Menge ist lokal zusammenhängend.
\end{vermutung}

\begin{vermutung}[Dichtheit der Hyperbolizität]
\label{verm:DHyp}
Hyperbolische Parameter sind dicht in $\mathcal{M}$.
\end{vermutung}

\begin{bemerkung}[MLC am Feigenbaum-Punkt — Gelöst]
Dudko und Lyubich~\cite{DudkoLyubich2023} haben bewiesen, dass
$\mathcal{M}$ beim Feigenbaum-Parameter
$c_{\mathrm{Feig}} \approx -1{,}4011552$ lokal zusammenhängend ist
(Satz~\ref{satz:DudkoLyubich}). Der allgemeine Fall für unendlich
renormalisierbare Parameter mit unbeschränkter Kombinatorik bleibt offen.
\end{bemerkung}

\begin{vermutung}[Siegel-Scheiben auf $\partial\mathcal{M}$]
Für jeden Parameter $c$ auf dem Rand einer Siegel-Scheiben-Komponente
ist $J_c$ eine Jordan-Kurve.
\end{vermutung}

\begin{bemerkung}
Vermutung~\ref{verm:MLC2} impliziert Vermutung~\ref{verm:DHyp}, aber
die Umkehrung ist nicht bekannt. Avila-Lyubich-de Melo haben wesentliche
Fortschritte bei der Renormalisierung erzielt, doch MLC bei unendlich
renormalisierbaren Punkten unbeschränkten Typs bleibt das zentrale
offene Problem der holomorphen Dynamik.
\end{bemerkung}

%% -------------------------------------------------------------------------
\section{Zusammenfassung}
%% -------------------------------------------------------------------------

Die Mandelbrot-Menge ist gleichzeitig ein Parameterraum, eine kombinatorische
Enzyklopädie aller quadratischen Dynamiken und ein konkretes Objekt, das die
tiefsten Merkmale holomorpher Iteration offenbart. Ihre Zusammenhängigkeit
(Douady-Hubbard) ist vollständig bewiesen; ihre lokale Zusammenhängigkeit
(MLC) ist das große offene Problem des Gebiets. Yoccoz' Puzzle-Technik löste
MLC für alle höchstens endlich-renormalisierbaren Parameter und reduzierte
das Problem auf den unendlich renormalisierbaren Fall. Die Äquivalenz mit der
Dichtheit der Hyperbolizität zeigt, dass MLC nicht bloß eine topologische
Aussage ist, sondern eine Strukturaussage über die gesamte quadratische Familie.

\begin{thebibliography}{99}

%% Alle Referenzen vollständig und korrekt

\bibitem{Mandelbrot1980}
B.~Mandelbrot,
\textit{Fractal aspects of the iteration of $z \mapsto \lambda z(1-z)$ for
complex $\lambda$ and $z$},
Ann.\ N.Y.\ Acad.\ Sci.\ \textbf{357} (1980), 249--259.

\bibitem{DouadyHubbard1982}
A.~Douady, J.~H.~Hubbard,
\textit{Itération des polynômes quadratiques complexes},
C.\ R.\ Acad.\ Sci.\ Paris \textbf{294} (1982), 123--126.

\bibitem{DouadyHubbard1984}
A.~Douady, J.~H.~Hubbard,
\textit{Étude dynamique des polynômes complexes, I, II},
Publications Mathématiques d'Orsay \textbf{84-02}, \textbf{85-04},
Université de Paris-Sud, Orsay, 1984--1985.

\bibitem{Thurston1985}
W.~P.~Thurston,
\textit{On the combinatorics and dynamics of iterated rational maps},
Preprint, Princeton University, 1985.

\bibitem{Milnor2000}
J.~Milnor,
\textit{Local connectivity of Julia sets: expository lectures},
in: \textit{The Mandelbrot Set, Theme and Variations},
London Math.\ Soc.\ Lecture Note Ser.\ \textbf{274},
Cambridge Univ.\ Press, 2000, 67--116.

\bibitem{HubbardSchleicher1994}
J.~H.~Hubbard, D.~Schleicher,
\textit{The spider algorithm},
in: \textit{Complex Dynamical Systems},
Proc.\ Sympos.\ Appl.\ Math.\ \textbf{49},
Amer.\ Math.\ Soc., 1994, 155--180.

\bibitem{Lyubich1997}
M.~Lyubich,
\textit{Dynamics of quadratic polynomials, I--II},
Acta Math.\ \textbf{178} (1997), 185--297.

\bibitem{KahnLyubich2009}
J.~Kahn, M.~Lyubich,
\textit{The quasi-additivity law in conformal geometry},
Ann.\ of Math.\ \textbf{169} (2009), 561--593.

\bibitem{Shishikura1987}
M.~Shishikura,
\textit{On the quasiconformal surgery of rational functions},
Ann.\ Sci.\ École Norm.\ Sup.\ \textbf{20} (1987), 1--29.

\bibitem{Shishikura1998}
M.~Shishikura,
\textit{The Hausdorff dimension of the boundary of the Mandelbrot set and
Julia sets},
Ann.\ of Math.\ \textbf{147} (1998), 225--267.

\bibitem{LavaursRousseau1993}
P.~Lavaurs,
\textit{Systèmes dynamiques holomorphes: explosion de points périodiques
paraboliques},
Thèse, Université Paris-Sud, Orsay, 1989.

\bibitem{DudkoLyubich2023}
D.~Dudko und M.~Lyubich,
\textit{Local connectivity of the Mandelbrot set at the Feigenbaum point},
arXiv:2309.02107, 2023.

\bibitem{AvilaLyubich2008}
A.~Avila, M.~Lyubich,
\textit{Feigenbaum Julia sets of positive area},
Preprint arXiv:math/0602458, 2006.

\end{thebibliography}

\end{document}
